\documentclass[10pt]{article}
\usepackage[margin=2.2cm]{geometry}
\usepackage[T1]{fontenc}
\usepackage[utf8]{inputenc}
\usepackage[polish]{babel}
\usepackage{amsmath,mathtools}
\usepackage{setspace}
\usepackage{booktabs,array}
\usepackage{enumitem}
\usepackage{hyperref}
\hypersetup{colorlinks=true,linkcolor=blue,urlcolor=blue,citecolor=blue}
\onehalfspacing
\setlength{\parindent}{1.2em}
\setlength{\parskip}{2pt}
\setlist[enumerate]{leftmargin=*,itemsep=4pt,topsep=3pt}

\title{Szybkie grupowanie opisów binarnych metodą CLASSIX\\
\large Protokół badawczy i opis projektu}
\author{Ivan Kovalov\\
Politechnika Opolska, Wydział Informatyki\\
Przedmiot: Rozpoznawanie wzorców, grupa 2\\
Projekt indywidualny}
\date{5 listopada 2026 r.}

\begin{document}
\maketitle

\begin{abstract}
Celem projektu jest zbadanie, jak liczba obiektów oraz stopień rozmycia grup wpływają na koszt obliczeniowy i jakość grupowania binarnych opisów danych. Główną analizowaną metodą jest CLASSIX w wariancie wykorzystującym odległość Tanimoto (CLASSIX\_T), porównywany z algorytmem DBSCAN przy tej samej definicji odległości. Badanie zostanie przeprowadzone na syntetycznych danych binarnych generowanych na podstawie wektorów bazowych i losowego odwracania bitów. Obejmie pomniejszoną adaptację eksperymentu z sekcji 5.1 publikacji Güttela i Roya oraz eksperyment własny ze zmiennym prawdopodobieństwem odwracania bitów. Jakość grupowania będzie oceniana za pomocą Adjusted Rand Index (ARI), liczby wykrytych klastrów i udziału szumu, a koszt -- przede wszystkim przez czas wykonania.
\end{abstract}

\section{Dane projektu}
\begin{singlespace}
\begin{tabular}{>{\bfseries}p{0.29\textwidth}p{0.65\textwidth}}
Temat & Szybkie grupowanie opisów binarnych metodą CLASSIX \\
Autor & Ivan Kovalov \\
Współautorzy & Brak -- projekt indywidualny \\
Uczelnia i wydział & Politechnika Opolska, Wydział Informatyki \\
Prowadzący & mgr inż. Michał Wierzbicki \\
Przedmiot / grupa & Rozpoznawanie wzorców, grupa 2 \\
Repozytorium projektu & \url{https://github.com/kovalov89/rozpoznawanie_wzorow26-27} \\
\end{tabular}
\end{singlespace}

\section{Cel i opis zadania}
Pytanie badawcze brzmi: \emph{Jak zmieniają się koszt i jakość grupowania opisów binarnych wraz z liczbą obiektów oraz rozmyciem grup?} Zadanie obejmuje wyjaśnienie reprezentacji binarnej, odległości Tanimoto oraz dwóch faz CLASSIX: agregacji obiektów wokół punktów początkowych i łączenia utworzonych grup. Należy również przedstawić regułę wczesnego odrzucania kandydatów właściwą dla odległości Tanimoto, sprawdzić dostępność implementacji i przygotować mały zbiór demonstracyjny. Zbiór demonstracyjny obejmuje 3 klastry, 16 cech i 18 obiektów (wektor bazowy oraz 5 obiektów pochodnych na klaster), z ziarnem generatora 42. Jest to pomniejszona demonstracja sposobu generowania danych, a nie pełne odtworzenie eksperymentów z publikacji.

\pagebreak
\section{Protokół badawczy}
\begin{singlespace}\small
\begin{enumerate}
  \item \textbf{Generowanie danych do eksperymentu głównego.} Zostaną utworzone 10 klastrów syntetycznych opisanych wektorami binarnymi o wymiarze $d=256$. Dla każdego klastra generator losuje wektor bazowy $x_i$. Następnie tworzy obiekty pochodne przez niezależne odwracanie każdego bitu z prawdopodobieństwem
  \[
  p_i=0{,}1\cdot\frac{\|x_i\|_1}{d},
  \]
  gdzie $\|x_i\|_1$ jest liczbą jedynek w wektorze bazowym. Przygotowane zostaną zbiory o liczebności 400, 800 i 1600 obiektów. Dla zachowania dokładnej liczebności każdy klaster będzie zawierał odpowiednio 40, 80 lub 160 obiektów łącznie z wektorem bazowym. Wektory bazowe, prawdopodobieństwa $p_i$, etykiety klastrów i ziarna losowe zostaną zapisane. Względem publikacji zmniejszono liczebność danych i wymiar do 256 cech; jest to jawnie oznaczona adaptacja generatora i eksperymentu.
  \item \textbf{Eksperyment własny -- rozmycie grup.} Dla $n=800$ i $d=256$ zostaną przygotowane trzy poziomy odwracania bitów, określone przez mnożniki $\lambda\in\{0{,}5;1;2\}$ bazowego prawdopodobieństwa: $p_i^{(\lambda)}=\lambda p_i$. Większy mnożnik oznacza większe rozproszenie obiektów wokół wektorów bazowych. Poziomy te są własnym rozszerzeniem badania.
  \item \textbf{Porównywane metody i parametry.} CLASSIX\_T zostanie porównany z DBSCAN. Obie metody otrzymają identyczne dane i będą korzystać z odległości Tanimoto. Jako punkt wyjścia przyjęte zostaną wartości z sekcji 5.1 publikacji: $\texttt{radius}=0{,}4$, $\texttt{minPts}=5$ dla CLASSIX\_T oraz $\texttt{eps}=0{,}4$, $\texttt{min\_samples}=5$ dla DBSCAN. Przed właściwymi pomiarami zostanie sprawdzona zgodność znaczenia parametrów z użytymi implementacjami; parametry zostaną następnie utrwalone w protokole.
  \item \textbf{Metryki jakości i kosztu.} Dla każdego uruchomienia zostaną zapisane: czas wykonania, Adjusted Rand Index (ARI) względem etykiet generatora, liczba wykrytych klastrów oraz udział obiektów oznaczonych jako szum. Etykiety generatora posłużą wyłącznie do oceny jakości, a nie jako wejście algorytmów.
  \item \textbf{Powtarzalność i kontrola pomiaru.} Zaplanowano pięć losowań z ziarnami 42--46. Obie metody będą testowane na tych samych realizacjach danych i w tym samym środowisku CPU. Zostaną zachowane parametry, wersje bibliotek, wektory bazowe, wyniki wszystkich uruchomień oraz błędy i przyczyny brakujących wyników. Wyniki zostaną przedstawione jako średnia i odchylenie standardowe. Jeżeli DBSCAN będzie wymagał wcześniejszego przygotowania macierzy odległości, czas jej obliczenia zostanie wliczony do czasu całkowitego; dodatkowo można podać sam czas grupowania. W porównaniu należy oddzielić wpływ algorytmu od różnic implementacyjnych. Dla wektorów zerowych zostanie jawnie przyjęta konwencja $d_T(0,0)=0$ oraz $d_T(0,b)=1$ dla $b\ne0$; zgodność tej konwencji w obu implementacjach zostanie sprawdzona.
\end{enumerate}
\end{singlespace}

\pagebreak
\section{Repozytorium i harmonogram checkpointów}
\begin{singlespace}
\textbf{Repozytorium projektu:}\\
\url{https://github.com/kovalov89/rozpoznawanie_wzorow26-27}\\[3pt]
\textbf{Repozytorium wykorzystanej implementacji CLASSIX:}\\
\url{https://github.com/kaustubhroy1995/classix_tanimoto_manhattan}

\begin{center}
\begin{tabular}{p{0.12\textwidth}p{0.60\textwidth}p{0.20\textwidth}}
\toprule
\textbf{Etap} & \textbf{Zakres checkpointu} & \textbf{Termin} \\
\midrule
C1 & Prezentacja tematu, protokół badawczy i mały zbiór demonstracyjny & 30.11.2026 \\
C2 & Generator oraz przygotowanie zbiorów 400, 800 i 1600 obiektów; test implementacji & 17.12.2026 \\
C3 & Porównanie CLASSIX\_T i DBSCAN, analiza metryk i wnioski z eksperymentów & 14.01.2027 \\
\bottomrule
\end{tabular}
\end{center}
\end{singlespace}

\begin{thebibliography}{9}
\bibitem{guettelroy2026}
S. Güttel, K. Roy, \emph{Fast and explainable clustering in the Manhattan and Tanimoto distance}, arXiv:2601.08781, 2026. Dostęp: \url{https://arxiv.org/abs/2601.08781}.

\bibitem{classixrepo}
K. Roy, \emph{classix\_tanimoto\_manhattan} [repozytorium kodu], GitHub. Dostęp: \url{https://github.com/kaustubhroy1995/classix_tanimoto_manhattan}.
\end{thebibliography}

\end{document}
